\section{Labeled multiscale completion and the erasure gate}
\label{sec:tpc63-labeled-completion}

The single-band obstruction does not force one to sum different cofactor
scales coherently.  There is a canonical finite construction which first
retains the scale as an orthogonal label.  This construction is exact at
the raw-energy level.  It is important, however, not to identify it with
the later physical output, where that label is absent.

\subsection{The canonical dyadic direct sum}

Let \(\Omega_X\) be the coefficient-independent finite literal incidence
which is being completed.  Thus an atom \(\omega\in\Omega_X\) carries, in
particular, its recovered cofactor \(r(\omega)\), and its pair projection
is the declared canonical relation (in the maximal eligible application,
this is \(\cI_X^{\max}\)).  No coefficient-dependent zero is removed from
\(\Omega_X\).

For \(k\in\Z\), set
\begin{equation}
 B_k=[2^k,2^{k+1})\cap\N,
 \qquad
 \Omega_{X,k}:=\{\omega\in\Omega_X:r(\omega)\in B_k\},
 \label{eq:tpc63-dyadic-atoms}
\end{equation}
and retain only the nonempty labels
\begin{equation}
 K_X^{\rm dyad}:=\{k:\Omega_{X,k}\ne\varnothing\},
 \qquad L_X:=|K_X^{\rm dyad}|.
 \label{eq:tpc63-nonempty-block-labels}
\end{equation}
If \(1\le r(\omega)\le X^{C+o(1)}\) for one fixed \(C\), as in the
inherited polynomial cofactor ranges, then
\begin{equation}
 L_X\le 2+\log_2\!\left(
       \frac{\max_{\omega\in\Omega_X}r(\omega)}
            {\min_{\omega\in\Omega_X}r(\omega)}
                         \right)
 =O(\log X)=X^{o(1)}.
 \label{eq:tpc63-logarithmic-label-count}
\end{equation}
Empty \(\Omega_X\) is harmless and may be treated separately; below we
assume it is nonempty.

Define
\begin{equation}
 \cH_X^{\rm raw}:=\ell^2(\Omega_X),
 \qquad
 \cH_{X,k}:=\ell^2(\Omega_{X,k}),
 \qquad
 \cH_X^{\rm lab}:=
       \bigoplus_{k\in K_X^{\rm dyad}}^{\perp}\cH_{X,k},
 \label{eq:tpc63-labelled-raw-spaces}
\end{equation}
all with raw counting-measure atomic norms.  The restriction map is
\begin{equation}
 \mathsf J_X:\cH_X^{\rm raw}\longrightarrow\cH_X^{\rm lab},
 \qquad
 \mathsf J_Xx=(x|_{\Omega_{X,k}})_{k\in K_X^{\rm dyad}}.
 \label{eq:tpc63-canonical-block-analysis}
\end{equation}

\begin{theorem}[Exact labeled multiscale completion]
\label{thm:tpc63-exact-labelled-completion}
The map \(\mathsf J_X\) is unitary and, for every raw atom vector \(x\),
\begin{equation}
 \boxed{
 \|\mathsf J_Xx\|_{\cH_X^{\rm lab}}^2
 =\sum_{k\in K_X^{\rm dyad}}
      \|x|_{\Omega_{X,k}}\|_2^2
 =\|x\|_{\cH_X^{\rm raw}}^2.}
 \label{eq:tpc63-exact-block-energy}
\end{equation}
At the pair level, writing
\(
 \cI_{X,k}^{\max}
 :=\{(m,r)\in\cI_X^{\max}:r\in B_k\}
\), one likewise has the disjoint identity
\begin{equation}
 \cI_X^{\max}
 =\coprod_{k\in K_X^{\rm dyad}}\cI_{X,k}^{\max}.
 \label{eq:tpc63-pair-block-completion}
\end{equation}
Consequently the passage to labeled dyadic blocks has exactly zero raw
energy loss and zero fixed-power bookkeeping cost.
\end{theorem}

\begin{proof}
The sets in \eqref{eq:tpc63-dyadic-atoms} are pairwise disjoint and their
union is \(\Omega_X\).  Parseval's identity for this coordinate partition
gives \eqref{eq:tpc63-exact-block-energy}; the inverse to \(\mathsf J_X\) simply
forgets the external block notation while retaining every atom.  The pair
identity follows from the same disjoint partition by the recovered
cofactor.  Finally, \eqref{eq:tpc63-logarithmic-label-count} implies
\(L_X^A=X^{o(1)}\) for every fixed \(A\), while the exact Hilbert identity
itself pays no factor at all.
\end{proof}

The theorem neither creates an arithmetically active atom nor asserts that
every block has a physical terminal realization.  It says only that once a
literal completed incidence has been declared, separating its cofactor
scales loses no raw mass.  In particular, the union in
\eqref{eq:tpc63-pair-block-completion} must not be substituted for a
coefficientwise physical identity that has been proved only on one of its
blocks.

\subsection{Exact analysis of label erasure}

We now isolate the additional map hidden by the phrase ``forget the
labels.''  Let \(\Xi_X\) be a finite output-coordinate set, let
\(\pi:\Omega_X\to\Xi_X\) be the declared grouping map, and attach a
nonzero complex weight \(a_\omega\) to every atom.  Define
\begin{equation}
 \cE_a:\ell^2(\Omega_X)\longrightarrow\ell^2(\Xi_X),
 \qquad
 (\cE_ax)(\xi)
 :=\sum_{\omega\in F_\xi}a_\omega x(\omega),
 \qquad
 F_\xi:=\pi^{-1}(\{\xi\}).
 \label{eq:tpc63-weighted-erasure-map}
\end{equation}
Coordinates with empty fiber carry no output information; henceforth they
are deleted, so every \(F_\xi\) is nonempty.  Put
\begin{equation}
 s_\xi:=\sum_{\omega\in F_\xi}|a_\omega|^2.
 \label{eq:tpc63-fiber-square-mass}
\end{equation}

\begin{theorem}[Fiber spectrum, norm, and kernel tests]
\label{thm:tpc63-erasure-fiber-tests}
For the grouping map in \eqref{eq:tpc63-weighted-erasure-map},
\begin{align}
 \cE_a\cE_a^*
 &=\operatorname{diag}(s_\xi)_{\xi\in\Xi_X},
 \label{eq:tpc63-erasure-co-Gram}\\
 \|\cE_a\|^2
 &=\max_{\xi:F_\xi\ne\varnothing}s_\xi,
 \label{eq:tpc63-erasure-exact-norm}\\
 \Ker\cE_a
 &=\left\{x:
       \sum_{\omega\in F_\xi}a_\omega x(\omega)=0
       \text{ for every }\xi\right\},
 \label{eq:tpc63-erasure-exact-kernel}\\
 \operatorname{rank}\cE_a
 &=\#\{\xi:F_\xi\ne\varnothing\},
 \qquad
 \dim\Ker\cE_a
 =\sum_\xi(|F_\xi|-1).
 \label{eq:tpc63-erasure-rank-nullity}
\end{align}
The nonzero squared singular values are exactly the numbers \(s_\xi\).
In particular, the full-domain lower constant
\begin{equation}
 \eta_a:=\inf_{x\ne0}
            \frac{\|\cE_ax\|_2^2}{\|x\|_2^2}
 \label{eq:tpc63-erasure-full-lower-constant}
\end{equation}
is positive if and only if every nonempty fiber is a singleton.  In that
case
\begin{equation}
 \eta_a=\min_{\omega\in\Omega_X}|a_\omega|^2;
 \label{eq:tpc63-erasure-singleton-lower}
\end{equation}
if some fiber contains at least two atoms, then \(\eta_a=0\).
\end{theorem}

\begin{proof}
The adjoint is
\((\cE_a^*y)(\omega)=\overline{a_\omega}y(\pi(\omega))\).
Hence \((\cE_a\cE_a^*y)(\xi)=s_\xi y(\xi)\), proving
\eqref{eq:tpc63-erasure-co-Gram} and
\eqref{eq:tpc63-erasure-exact-norm}.  The displayed fiber equations give
the kernel.  A nonempty fiber contributes one to the rank and
\(|F_\xi|-1\) to the nullity, proving
\eqref{eq:tpc63-erasure-rank-nullity}.  The positive spectrum of
\(\cE_a^*\cE_a\) equals that of \(\cE_a\cE_a^*\).  If a fiber contains
two atoms, a nonzero weighted difference supported on those atoms lies in
the kernel.  If every fiber is a singleton, the Gram is diagonal with
entries \(|a_\omega|^2\), which proves
\eqref{eq:tpc63-erasure-singleton-lower}.
\end{proof}

For the literal unweighted erasure \(a_\omega=1\), the exact formulas
specialize to
\begin{equation}
 \|\cE_1\|^2=\max_\xi|F_\xi|,
 \qquad
 \Ker\cE_1
 =\left\{x:\sum_{\omega\in F_\xi}x(\omega)=0
                   \text{ for every }\xi\right\}.
 \label{eq:tpc63-unweighted-erasure-formulas}
\end{equation}
If the only discarded coordinate is the dyadic block label and each block
contributes at most one atom to a given output coordinate, then
\begin{equation}
 \|\cE_1\|^2\le L_X=O(\log X)=X^{o(1)}.
 \label{eq:tpc63-pure-block-erasure-upper}
\end{equation}
This is an upper norm, not a lower reassembly estimate.  Two blocks which
land on the same physical coordinate already give the kernel vector
\((1,-1)\) on that fiber.

Normalizing each fiber by \(|F_\xi|^{-1/2}\) would make the nonzero
singular values equal to one, but it would neither remove the fiber kernel
nor preserve the literal raw atomic normalization.  Such a normalization
therefore cannot be silently inserted into the physical packet.

\subsection{The relevant test is on the actual upstream range}

The full labeled space may be larger than the image of the actual
coefficient map.  Let
\(\mathsf A_X^{\rm lab}:V_X^{\rm coef}\to\cH_X^{\rm lab}\) be any declared
upstream analysis map and put \(S_X=\Ran\mathsf A_X^{\rm lab}\).  Assume
\(S_X\ne\{0\}\), and define
\begin{equation}
 \eta(\cE_a;S_X)
 :=\inf_{0\ne v\in S_X}
       \frac{\|\cE_av\|_2^2}{\|v\|_2^2}.
 \label{eq:tpc63-range-erasure-constant}
\end{equation}

\begin{proposition}[Exact range and composition criteria]
\label{prop:tpc63-range-erasure-criterion}
In finite dimension,
\begin{equation}
 \eta(\cE_a;S_X)>0
 \quad\Longleftrightarrow\quad
 S_X\cap\Ker\cE_a=\{0\},
 \label{eq:tpc63-range-kernel-iff}
\end{equation}
and \(\eta(\cE_a;S_X)\) is the least eigenvalue of the compression of
\(\cE_a^*\cE_a\) to \(S_X\).  If
\begin{equation}
 \|\mathsf A_X^{\rm lab}z\|_2^2\ge\alpha_X\|z\|_2^2,
 \qquad
 \eta(\cE_a;S_X)\ge\eta_X,
 \label{eq:tpc63-two-composable-lowers}
\end{equation}
then
\begin{equation}
 \|\cE_a\mathsf A_X^{\rm lab}z\|_2^2
 \ge\eta_X\alpha_X\|z\|_2^2.
 \label{eq:tpc63-composed-erasure-lower}
\end{equation}
Conversely, a nonzero \(z\) with
\(\mathsf A_X^{\rm lab}z\in\Ker\cE_a\) kills every full-coefficient lower bound for this
particular composition.
\end{proposition}

\begin{proof}
The unit sphere of \(S_X\) is compact.  The continuous function
\(v\mapsto\|\cE_av\|^2\) has positive minimum exactly when it has no zero
there, which proves \eqref{eq:tpc63-range-kernel-iff}.  Rayleigh--Ritz gives
the compressed-Gram statement.  Applying the two inequalities in
\eqref{eq:tpc63-two-composable-lowers} successively proves
\eqref{eq:tpc63-composed-erasure-lower}; the final assertion is immediate.
\end{proof}

Thus the honest implication is
\begin{equation}
 \boxed{
 \text{exact labeled raw completion}
 \quad\not\Longrightarrow\quad
 \text{physical grouped lower bound}.}
 \label{eq:tpc63-labelled-not-physical}
\end{equation}
One must first identify the literal grouping map and the actual upstream
range, apply the kernel test, and only then estimate the compressed Gram.
The logarithmic number of labels controls bookkeeping and certain upper
norms; it does not control coherent cancellation after grouping.
